\chapter{Background and Related Work}\label{ch:background}
% Lecture 3: the review is the background of the research questions, not a list of papers.
% Each section ends with what is missing for this project.
% Draft v0.1, 10 November 2026: selective reading of nine works (search log in the reference
% manager). Sections marked TODO wait for works not yet read.

This chapter places \project{} between four lines of work: understanding visually rich
documents, checking documents for completeness and consistency, generating answers with
citations, and detecting changes on web pages. Each section ends with what the work leaves open
for a visa application file. The review is selective, not systematic: nine works were read in
the sections relevant to the research questions, and the search terms and databases are logged
in the reference manager.

\section{The Schengen short-stay visa in Türkiye}
The supporting documents for applicants in Türkiye are harmonised by the European Commission
\autocite{eu2021harmonised}; each member state adds national requirements, for example the
Netherlands \autocite{nl2025checklist} and Norway \autocite{no2025checklist}.
Table~\ref{tab:nl-no} shows that the two national layers differ in exactly the kind of detail an
applicant can miss: a translation requirement, a general limit on document age and the way
passport validity is counted.

\begin{table}[ht]
  \centering
  \caption{National differences on top of the EU baseline (sources reviewed on 23 October 2026).}
  \label{tab:nl-no}
  \begin{tabularx}{\textwidth}{@{}lXX@{}}
    \toprule
    Topic & Netherlands & Norway \\
    \midrule
    Translation & not required & all documents in English or Norwegian \\
    General document age & not stated (photo $\leq$ 6 months) & no document older than 6 months \\
    Passport validity & 3 months after leaving the Schengen area & 3 months beyond the intended stay \\
    Financial means & bank statements of 3 months or proof of income & salary slips and bank movements of 3 months \\
    \bottomrule
  \end{tabularx}
\end{table}

The official pages also change. In the snapshot of 5 November 2026, three of the five source
pages differed from their first snapshot of 23 October; on inspection, two of them had no visible
change and one had gained a link to a letter of authorisation (Section~\ref{sec:change-detection}).

\section{Document understanding and information extraction}
Extracting fields from scans and photos of documents is a mature research area.
\textcite{ding2026mllm} survey methods built on multimodal large language models and separate
OCR-based pipelines, which read the text first and then reason over it, from OCR-free models,
which map the image directly to structured output. They list the recurring problems: errors of
the OCR step carried into later steps, complex layouts, links across pages and training data
dominated by English.

Donut \autocite{kim2022donut} is a representative OCR-free model: an image encoder and a text
decoder turn a document image into JSON. It reaches a field-level F1 of 84.1 on receipts (CORD)
but does not beat OCR-based models on every task (DocVQA ANLS 67.5 against 78.1 for LayoutLMv2),
and no Turkish documents are tested. Its synthetic generator, SynthDoG, renders text on varied
backgrounds with noise, blur and perspective; it creates visual variety, not the missing or
contradictory documents that \project{} has to detect.
\todo{Kazamel, Kocasarı and Alıcı (2025), Turkish ID cards: add once the full text is read.}

\paragraph{What is open.} These works measure extraction from one document at a time. Whether
an application file as a whole is complete and consistent with the requirements is not their
question, so this project reports field extraction and error finding as separate measures (AS1).

\section{Checking documents for completeness and consistency}
The closest work is DocChecker \autocite{chen2026contradictions}. It builds entity--relation
graphs from text and finds seven kinds of contradiction within and across documents, combining
logic rules with machine-learned predicates; the rules are learned with the help of a language
model, and external knowledge graphs can be attached. It reaches an average F1 of about 0.90 on
15 text benchmarks, none of which are application files. A hybrid of rules and learning is
therefore not new.

Nor is validation after extraction. \textcite{bulatov2022identity} describe a framework for
identity documents read from scans, photos and video, in which field syntax and relations between
fields (an issue date cannot precede the date of birth) are kept in configuration and the result
is checked for completeness. The rules are configured by hand, and the paper reports no new
end-to-end experiment.

Whether a language model alone can do this work is doubtful. On PRISMM-Bench
\autocite{selch2025prismm}, 384 inconsistencies reported by reviewers in scientific papers, 21
multimodal models reach an accuracy of 27.8 to 53.9\,\% in multiple-choice tasks where chance is
25\,\%. This supports including a model-only baseline in AS1; it does not show that a hybrid
system does better, which AS1 has to measure on this project's own error types.

Comparing names across documents is a problem of its own. \textcite{christen2006names} finds
that no single matching technique is best: the ranking changes with the part of the name, word
order and the kind of error. The name rule of this project therefore keeps the original name,
normalises Turkish letters and spacing, and chooses its threshold on labelled pairs.

\paragraph{What is open.} In these works the rules are configured or learned for one domain; none
of them takes requirements from official, dated sources with the passage cited, checks that every
required document of a multi-document file is present, or follows a change of those requirements
to the files it affects.

\section{Retrieval-augmented generation and citation}
An assistant that answers questions about requirements has to show where each statement comes
from. ALCE \autocite{gao2023alce} measures this: a sentence counts as supported when its cited
passages entail it (citation recall), and a citation counts as precise when it is needed for that
support (citation precision), both judged by an inference model. Even strong models leave many
sentences unsupported; on long-form questions (ELI5), GPT-4 with 20 passages reaches a citation
recall of 48.5\,\% and a precision of 53.4\,\%.
\todo{Survey of retrieval-augmented generation \autocite{gao2024ragsurvey} and of hallucination
\autocite{ji2023hallucination}: add once read.}

\paragraph{What is open.} ALCE asks whether a passage supports a sentence, not whether the source
is official, current or the version in force. This project adds the second part by tying each
requirement to a versioned requirement set built from dated snapshots of official pages.

\section{Change detection of web pages}\label{sec:change-detection}
\textcite{mallawaarachchi2020change} survey change detection and notification: kinds of change
(content, structure, presentation, behaviour), fixed and adaptive check schedules, and filtering
of irrelevant parts such as advertisements. There is no shared benchmark. \textcite{adar2009web}
crawled about 55,000 frequently revisited pages every hour for five weeks; 11.5\,\% of them
already differed between two simultaneous requests, because of advertisements, timestamps and
counters. A changed hash is therefore only a candidate change, which matches the first
observation on the official visa pages above.
\todo{Semantic change detection \autocite{zhang2004xml}: add once read.}

\paragraph{What is open.} Deciding whether a change affects a visa requirement needs knowledge of
the requirements. AS3 labels pairs of snapshots as meaningful, minor or cosmetic and measures
precision, recall and F1 per class; small edits such as a changed number of months can matter,
so changes are not filtered by size alone.

\section{Similar systems}
A search of public repositories on 23 October 2026 found several open-source visa and
document-check assistants (Table~\ref{tab:similar}). None of the repositories inspected reports a
measured error rate.

\begin{table}[ht]
  \centering
  \small
  \caption{Similar open-source systems (repositories inspected on 23 October 2026).}
  \label{tab:similar}
  \begin{tabularx}{\textwidth}{@{}lXX@{}}
    \toprule
    System & Approach & Difference from \project{} \\
    \midrule
    VISAPATH &
      Completeness and consistency checks of visa documents with retrieval and a language
      model as judge; scans uploads for prompt injection &
      Sources not named; no explicit, tested rules \\
    UK Visa Auditor &
      JSON rule engine, OCR and fuzzy matching for UK student visas &
      Links to official pages in chat only; no versioning \\
    Visa-Agent &
      Model proposes facts, Python rules decide; versioned rule snapshots &
      Same principle; sources and change detection not documented \\
    changedetection.io &
      General web page monitor with selector filters, text diff and notifications &
      Detects changes but cannot tell which requirements or applicants they affect \\
    \bottomrule
  \end{tabularx}

  \raggedright\footnotesize Repositories: \url{https://github.com/sankeerthanayerubandi-gif/VISAPATH}; \url{https://github.com/taimoorcodes-cmd/uk-visa-auditor}; \url{https://github.com/Sanssssssssssssssss/Visa-Agent}; \url{https://github.com/dgtlmoon/changedetection.io}.
\end{table}

\section{Research gap}
Hybrid checking with rules and learned components exists \autocite{chen2026contradictions}, and
so does validation of extracted fields \autocite{bulatov2022identity}. What the works reviewed
here do not show together is a check of a multi-document visa application against requirements
taken from official, dated and versioned sources with the passage cited, which reports missing
documents as well as inconsistent fields, and which carries an approved change of a source to the
applications it affects. Whether such a hybrid finds errors better than a language model alone
(AS1), how it compares with human reviewers (AS2), and how well source monitoring separates
meaningful from cosmetic changes (AS3) are the questions this project measures. The claim is
limited to the works reviewed; it is not the result of a systematic review.
